\documentclass[a4paper,11pt,oneside]{book}

\usepackage[english]{babel}
\usepackage[showframe=false]{geometry}
\usepackage[usenames,dvipsnames]{xcolor}
\usepackage[utf8]{inputenc}
\usepackage[T1]{fontenc}
\usepackage{changepage}
\usepackage[]{algorithm2e}
\usepackage{amssymb}
\usepackage{amsmath}
\usepackage{graphicx}
\usepackage{listings}
\usepackage{verbatimbox}
\usepackage{ulem}
\usepackage{fancyvrb}
\usepackage{float}
\usepackage{hyperref}
\usepackage[parfill]{parskip}
\usepackage{tikz}
\usepackage{pdflscape}
\usepackage{minted}
\usepackage{titlesec}
\usepackage{titleps}
\usepackage{lastpage}
\usepackage{fancyhdr}
\usepackage{etoolbox}
\usepackage[]{algorithm2e}
%Following overwrites the page style for chapters
\patchcmd{\chapter}{\thispagestyle{plain}}{\thispagestyle{ruledChapter}}{}{}

%New page style for chapters
\newpagestyle{ruledChapter}{
	\setfoot{}{\thepage\ of \pageref{LastPage}}{}
	\footrule
	\renewcommand\makefootrule{\color{black}\rule[\baselineskip]{\linewidth}{0.4pt}}
}
%New page style for rest
\newpagestyle{ruled}{
	\sethead{\raggedright \chaptername\ \thechapter :\ \chaptertitle}{}{}
	\headrule
	\setfoot{}{\thepage\ of \pageref{LastPage}}{}
	\footrule
	\renewcommand\makeheadrule{\color{black}\rule[-.3\baselineskip]{\linewidth}{0.4pt}}
	\renewcommand\makefootrule{\color{black}\rule[\baselineskip]{\linewidth}{0.4pt}}
}

\expandafter\def\csname PY@tok@err\endcsname{}
\newcommand{\HRule}{\rule{\linewidth}{0.5mm}}
\newcommand{\specialcell}[2][c]{%
  \begin{tabular}[#1]{@{}c@{}}#2\end{tabular}}

\addtocontents{toc}{\protect\thispagestyle{empty}}

\title{}
\author{}
\date{} 

\begin{document}
\begin{titlepage}
\begin{center}

%-----------------------------------------------------------------
%							FRONTPAGE
%-----------------------------------------------------------------
\thispagestyle{empty}
\includegraphics[width=0.55\textwidth]{logo.pdf}\\[1cm]    
\textsc{\Large DM818 Assignment 3}\\[0.5cm]

% Title
\begin{Huge}
\textbf{DNS Algorithm}
\end{Huge}

\vspace{4cm}

% Author and supervisor
\begin{minipage}{1\textwidth}
\begin{center}
\emph{}\\

Dan \textsc{Sebastian Thrane}\\
\verb!<dathr12@student.sdu.dk>!\\

Lars \textsc{Thomasen}\\
\verb!<latho12@student.sdu.dk>!\\

\end{center}
\end{minipage}
\begin{minipage}{0.4\textwidth}
\end{minipage}

\vfill
% Bottom of the page
{\large Winter 2015}\\

\end{center}
\end{titlepage}

%-----------------------------------------------------------------
%							   TOC
%-----------------------------------------------------------------
\renewcommand{\contentsname}{Table of Contents}
\tableofcontents
\thispagestyle{empty}

%-----------------------------------------------------------------
%						  ACTUAL REPORT
%-----------------------------------------------------------------
\pagestyle{ruled}
\chapter{Introduction}
\setcounter{section}{1}
This report documents the development of an MPI virtual topology, and the implementation of the DNS algorithm covered in DM818.

The implementation takes some assumptions, such as it expects a cubic number $p$ of processors and assumes that the matrices of size $n*n \geq p^{1/3}$.

\section{Contributions}
Programming was done only using pair programming, and as such an equal contribution of this assignment has been done.

\chapter{Implementation}
\section{Design choices}
%A description of the design choices that you tried and how they affected the performance.
No specific design choices has been made. The implementation consists of two phases:
\begin{itemize}
\item Create virtual topology using MPI.
\item Follow the DNS algorithm.
\end{itemize}

Creating the virtual topology is a rather simple task, when a good understanding of how MPI works, and thus little to no design choices will be needed. 

Implementation of the DNS algorithm left no room for design choices, as it was simply a matter of following the algorithm presented in the course.

\section{Virtual topologies}
%A description of how you used virtual topologies.
The communication pattern of a set of processes can be represented by a graph. The nodes in this graph is represented by processes and the edges between them is the processes communication between each other. MPI provides message-passing between any pair of processes in a group, often represented in a Communicator.

MPI offers built-in methods such as \verb!MPI_Graph_create! and \verb!MPI_Cart_create!. The latter offers the possibility to describe Cartesian structures in more than two dimensions.

The virtual topology was created in MPI using \verb!MPI_Cart_create!, this method takes the following parameters:
\begin{itemize}
\item MPI Communicator to which the Cartesian topology information is attached.
\item Number of dimensions.
\item Number of processes for each dimension.
\item Whether the dimensions are periodic.
\item Whether ranking may be reordered, if false the rank from the old communicator is retained.
\end{itemize}
and outputs a new communicator for the Cartesian topology. This is used to create a 3-dimension hypercube.

In order to communicator between the processes in this hybercube, we can create additional Communicators using \verb!MPI_Cart_sub!, specifying which dimensions we want our new communicator to attach to processes in.

As an example the following line specifies that we want all processes in the i'th axis in the new communicator \verb!iComm!.

\begin{verbatim}
int iDimensions[3] = {1, 0, 0};
MPI_Cart_sub(gridCommunicator, iDimensions, &iComm);
\end{verbatim}

\section{Matrix not divisible by $p^{1/3}$}
PADDING! TODO THIS SECTION!

\chapter{Performance}
\section{Runtimes}
%In your report, all runtimes and measures that depend on runtimes, should be based on average 
%runtimes based on 10 test runs (compare main.c).
\chapter{Isoefficiency of the implementation}


\chapter{Conclusion}


%-----------------------------------------------------------------
%						     APPENDIX
%-----------------------------------------------------------------
%/newpage
%\newgeometry{left=2.5cm,right=2.5cm}
%\chapter{Appendix}
%\section{Serial algorithm plotting data}

\end{document}
